% Molt Petit — the from-scratch rewrite (2026-08).
% Ground rules for this document (agreed with George):
%  - written for the person who would DEPLOY the protocol; accurate, concise;
%  - target ~15 pages, no appendix unless a detail is truly needed but would
%    break the main flow;
%  - implementation + experiments get at most ONE page total, with one
%    concise measurements table; no slot-time formulas;
%  - every machine-checked claim cites its Lean name: a theorem `xxx` is
%    MoltPaper.xxx (statement thm_xxx in Spec/Statements.lean); definitions
%    live in Spec/. Formal detail belongs in the Lean, not in the prose.
\documentclass[10pt]{article}

\usepackage[a4paper,margin=2.6cm]{geometry}
\usepackage{amsmath,amssymb,amsthm}
\usepackage{booktabs}
\usepackage{microtype}
\usepackage[hidelinks]{hyperref}
\usepackage{xcolor}
\usepackage{graphicx}

\definecolor{yakink}{HTML}{3A3733}
\definecolor{yakaccent}{HTML}{C3401F}

\newtheorem{theorem}{Theorem}
\newtheorem{assumption}{Assumption}
% \code: typewriter identifiers; long Lean names may break (hyphenless)
% after an underscore.
\DeclareRobustCommand{\code}[1]{{\renewcommand{\_}{\textunderscore\allowbreak}\texttt{#1}}}
\newcommand{\quorum}{\ensuremath{q}}
\newcommand{\fmax}{\ensuremath{f}}

\title{\textbf{Molt Petit: A Machine-Checked Slot-Based BFT Consensus
Protocol with Recursive Certificates}}
\author{George Agapov\\ \normalsize Yak Software, Andorra}
\date{September 2026, draft}

\begin{document}

\noindent{\LARGE\bfseries\raggedright\hyphenpenalty=10000 \exhyphenpenalty=10000
Molt Petit: A Machine-Checked Slot-Based BFT Consensus
Protocol with Recursive Certificates\par}
\vspace{4pt}
\noindent\begin{minipage}[b]{\dimexpr\textwidth-26mm\relax}
  George Agapov \textperiodcentered\ Yak Software, Andorra\par
  \vspace{2pt}
  {\small\color{yakink}September 2026, draft \textperiodcentered\ yak.finance/molt-petit\par}
\end{minipage}%
\hfill
\begin{minipage}[b]{20mm}
  \hfill\includegraphics[width=\linewidth]{yak-wordmark.pdf}%
\end{minipage}
\vspace{1em}

\begin{abstract}
Molt Petit is a consensus protocol for a fixed set of $n$ validators
that take turns producing blocks in time slots. A light client needs
very little: it keeps the genesis block and a clock, downloads a small
certificate from any source, checks that the certificate verifies and
that its newest block is recent, and then trusts every block at least
$n$ blocks deep. A chain is valid only if, in every past window of $n$
slots, at least two thirds of the validators produced a block. There
are no voting rounds, and finality is at a fixed depth.

All safety claims are machine-checked in Lean~4, under the assumptions
listed in Section~\ref{sec:assumptions} and a few implementation
assumptions named in Section~\ref{sec:impl}. The main safety theorem
and its form for each rotation mode cover the Rust client that ships,
end to end: from the certificate it verifies, through the block ids and
recency it checks itself, to the blocks it accepts. The recursive
certificate runs a verified step per block, written in the same source
and proved to build exactly the claims the theorems assume; we trust the
proof system and the circuit's gadgets to run it faithfully.

The attacker may obtain validators' signing keys at any time --- by
theft or by a temporary break-in --- and sign anything with them. Safety
holds as long as few enough keys are exposed around each window's own
time. Checking that the certificate is recent is essential: it is what
makes keys stolen long ago useless, and the proofs depend on it. No
valid chain can claim to be far ahead of real time.

Validators can rotate keys inside the protocol. Every rotation mode
recovers from lost keys. Against an attacker who steals keys and keeps
them forever, two modes stay safe given clearly stated duties for
operators, and clients still need only the genesis, a clock and, in
one of the two, the public rotation schedule. With post-quantum signatures on ordinary hardware, the
verified consensus rule adds under $2\%$ to the cost of producing a
certificate.
\end{abstract}

\section{Introduction}
\label{sec:intro}

Think of a light client that is offline most of the time: a wallet, an
embedded device, a bridge contract. Now and then it has to act on the
current state of a small consensus system. Today it has three options:
replay the history, trust a service, or verify a chain of committee
hand-offs.

Molt Petit is built so that this client needs only three things:
\begin{itemize}
\item the deployment's genesis block,
\item a clock, and
\item a small, constant-size \emph{certificate}, fetched from anyone.
\end{itemize}
If the certificate verifies, at least one signed block follows it, and
the newest block is recent, the client can act on everything at least
$n$ blocks below that newest block. Two clients that follow this rule at
the same time never act on conflicting blocks; Section~\ref{sec:lc}
says when clients checking at different times are covered.

This guarantee is a machine-checked theorem about the validator code that
actually ships (\code{rust\_deployed\_agreement}, and
\code{ts\_deployed\_agreement} for the TypeScript version), under the
assumptions of Section~\ref{sec:assumptions}. That code checks the
certificate, the block ids and the recency itself.

\paragraph{A small protocol.} A fixed roster of $n$ validators takes turns
producing blocks, one per time slot, in a public order. There are no
votes and no view changes. A chain is valid only if every past window of
$n$ slots contains blocks from at least two thirds of the validators. We
call this the \emph{density} rule (Section~\ref{sec:protocol}).

Density makes deep forks impossible while the fault budget holds, so a
block $n$ deep is final; finality is a fixed depth, not a probability.
Fork choice matters only for liveness: a node simply follows the highest
valid tip. A node never stores history. It holds a recursive proof that
the chain so far is valid, plus a tail of signed blocks
(Section~\ref{sec:rules}).

\paragraph{Recent certificates matter.} Over a deployment's lifetime,
keys leak: machines are broken into and later cleaned, keys are stolen.
An attacker can use old leaked keys to sign a fake history; this is the
classic long-range attack. Molt Petit's fault budget counts, for each
window of $n$ slots, only the keys exposed during a short stretch of real
time around that window. A client that rejects certificates older than
$\varphi$ slots is exactly the client for whom later leaks do not matter.
And no attacker can prepare a valid chain that runs far ahead of real time to
get around this check (Section~\ref{sec:timed}).

\paragraph{Key rotation is part of consensus.} Every block says which
version of its producer's key signed it, so the chain orders key
rotations like any other event. Every rotation mode recovers from a lost
key, with no extra work for clients. Against an attacker who steals keys
and keeps them forever, two modes are safe: a public rotation schedule,
and a lockstep counter with key erasure. Each comes with a short, exact
list of operator duties, and their clients may sleep for as long as they
like. The third mode, reactive rotation, recovers from lost keys only; we
show why theft defeats it (Section~\ref{sec:rotation}).

\paragraph{Proofs about the shipped client.} The client is one Rust
file. We translate it into Lean~4 (with Charon and Aeneas) and prove the
translation sound against the model. The safety theorem, in every
rotation mode, therefore holds for the code that runs, from the
certificate it verifies to the blocks it accepts, with signature checking
modelled as an ideal functionality (Section~\ref{sec:assumptions}). The
other key-rotation results are proved for the model. The same source
holds the recursive certificate's step: per block, it checks links,
density and the mode's key rule, recomputes the block's id, and is proved
to build exactly the claims the theorems assume. We trust the proof
system and the circuit's gadgets to run it faithfully. A TypeScript
client is proved sound against the same model, which keeps the model
independent of any one implementation. No theorem assumes anything
beyond classical logic.

Measured with post-quantum signatures on ordinary hardware, the verified
consensus rule costs under $2\%$ of a certificate's proving time
(Section~\ref{sec:impl}).

\paragraph{The price.} The simplicity has costs, stated here and in
Section~\ref{sec:limitations}: the roster is closed and equally weighted;
a window without enough blocks halts the chain, by design; and safety
relies on correct clocks rather than holding under arbitrary network
delays. Where these trades are acceptable (permissioned systems, bridges,
fleets of devices), the result is a consensus layer whose client rule
fits in one sentence, and whose proofs re-check from a fresh checkout.

\section{How to run it}
\label{sec:opmodel}

This section says what operators and clients must do. The rest of the
paper makes it precise.

One protocol runs throughout. A fixed roster of $n$ validators produces
blocks in public round-robin slots. Every block states which version of
its producer's key signed it. A node holds a constant-size certificate
plus a tail of recent signed blocks. A verifier accepts a chain only if
it is valid, its tail holds at least one block, and its newest block is
recent, at most $b$ slots old on the verifier's own clock. It then
trusts everything $n$ blocks deep, or $2n$ deep in mode 3 under the
per-generation budget, counted in the history the certificate attests
followed by the tail (Section~\ref{sec:rotation}).

Key rotation is built in. A deployment picks one of three rotation
\emph{modes}. The modes differ in one rule, the \emph{pin}: the lowest
key version a block may use (modes 1 and 2), or the exact version for
each window of slots (mode 3). The pin decides when a retired key stops
being usable.

The three modes behave the same unless keys are \emph{stolen}.
Validators rotate keys; clients never react to rotations. Every client
duty below is a fixed rule in time, not a response to an event.

\paragraph{Lost keys: every mode is safe.}
Any attacker in this paper can make a validator \emph{lose} its key:
destroy it, jam its signer, or force a rotation. Against that, all three
modes give the same client contract: \textbf{keep the genesis and a
clock, and sleep as long as you like}. A destroyed key is held by no
one, so losing it exposes nothing, provided the validator's record of
what it has signed survives the loss; Theorem~\ref{thm:lc} then holds
for each mode's validator (Section~\ref{sec:rotation}). The validator that
lost its key recovers through its mode: a new key version in mode 1, the
next scheduled generation in mode 2 (offline for up to one era), the
next lockstep switch in mode 3. In mode 1, Lean proves that its next
block, stating the new version, extends a valid unsigned chain whose
windows the others kept dense (\code{liveness\_produce\_blockK}).

\paragraph{Mode 1: reactive rotation.}
Validators rotate whenever they choose, routinely or after losing a key.
Nobody knows in advance when a rotation will happen. Clients keep the
genesis and a clock and get the lost-key guarantee above, and nothing
more.

Mode 1 is \textbf{not safe against theft}. An attacker who steals the
retired keys of enough validators can forge a recent fork that branches
off before those rotations (Section~\ref{sec:rotation}). A deployment
that must survive theft runs mode 2 or mode 3.

\paragraph{Mode 2: scheduled rotation.}
The key version in force at each slot is a \emph{public function of the
slot}, published before the deployment starts. Everyone knows the whole
schedule in advance, and the validator rejects versions below it.

In every mode, signing keys are derived from one offline master key per
validator (Section~\ref{sec:protocol}). Mode 2 adds timing: the operator
derives each period's key (a \emph{generation}) just before that period
(its \emph{era}) begins. The theft theorem accounts for this through its
fault budget instead of assuming it (Section~\ref{sec:rotation}).

The client keeps the genesis, a clock, and the schedule. It may stay
offline as long as it likes and sync whenever it wakes up. There is no
trusted checkpoint, no required check-in rhythm, and no key erasure.

\paragraph{Mode 3: lockstep rotation.}
There is no schedule at all. The roster moves one shared generation
counter forward whenever it chooses. The slots are cut into fixed
windows of $n$ slots. A valid chain never mixes generations inside one
window, and the counter never goes back. At each switch, retired keys
are physically destroyed, in hardware security modules or RAM-only
storage; storage that does not truly erase does not qualify.

The client keeps the genesis and a clock, even less than in mode 2, and
may sleep as long as it likes. The price is paid by operators: they must
coordinate the switches, and erasure must be real. In return, erasure
earns a finer fault budget, counted per generation
(Section~\ref{sec:rotation}).

\paragraph{The modes at a glance.}
\begin{center}
\small
\setlength{\tabcolsep}{3.5pt}
\begin{tabular}{@{}p{1.6cm}p{2.3cm}p{2.5cm}p{1.4cm}p{6.2cm}@{}}
\toprule
 & pin comes from & client keeps & erasure & main Lean theorem \\
\midrule
no theft, any mode & any of the three &
  genesis + clock & no &
  Thm.~\ref{thm:lc}: \code{timed\_light\_client\_safety}
  (Rust: \code{rust\_timed\_certified\_agreement}); restricted to each
  mode's validator: \code{keyrot\_loss\_agreement},
  \code{sched\_loss\_agreement}, \code{lockstep\_loss\_agreement} \\
\addlinespace
1 reactive & the chain itself & --- & --- &
  none against theft (Section~\ref{sec:rotation}); lost keys: row above \\
\addlinespace
2 scheduled & the public schedule &
  genesis + clock + schedule & no &
  Thm.~\ref{thm:sched}: \code{sched\_exposure\_agreement},
  \code{sched\_exposure\_certified\_agreement} \\
\addlinespace
3 lockstep & the roster's counter &
  genesis + clock (clock only: Thm.~\ref{thm:lock-gen}, depth $2n$) & yes &
  Thm.~\ref{thm:lock-gen}: \code{lockstep\_client\_safety\_gen};
  Thm.~\ref{thm:lock} (theft budget is lifetime):
  \code{lockstep\_client\_safety} \\
\bottomrule
\end{tabular}
\end{center}
Reading down the table, each mode trusts something different. Mode 1
trusts that keys are only ever lost, never stolen. Mode 2 trusts the
public schedule, and charges key handling to its fault budget. Mode 3
trusts roster coordination and physical key destruction.

\section{The protocol}
\label{sec:protocol}

Lean names appear in \code{typewriter}. A cited theorem \code{xxx} is
proved as \code{MoltPaper.xxx} in the file \code{Paper/Proofs.lean}. Its
exact statement is \code{thm\_xxx}, in the file
\code{Spec/Statements.lean}. Unless stated otherwise, results assume a
non-empty roster, $n \ge 1$; the liveness and sanity statements, the
full-chain soundness statements of the Rust validator and the lockstep
package conversions need no roster hypothesis, and
\code{groundedCertLock\_gen\_of\_tail} needs $n \ge 2$. Every cited
definition is in \code{Spec/}: a few files that hold the definitions and
statements only, and that the rest of the development imports. The
precise forms live there; this paper explains them.

\subsection{Setting}

Time is divided into slots of fixed length $\tau$. A fixed, closed set of
$n$ validators, the \emph{roster}, produces blocks. Slot $s$ belongs to
validator $s \bmod n$ (\code{producer}). Who produces, and when, is
public and fixed; there is no leader election.

Two numbers derived from $n$ appear everywhere:
\begin{itemize}
\item the \emph{quorum} $\quorum = \lceil 2n/3 \rceil$, and
\item the \emph{fault budget} $\fmax = \lfloor (n-1)/3 \rfloor$: how many
  slots in any window of $n$ slots the attacker may control
  (Assumption~\ref{ass:budget}).
\end{itemize}
Table~\ref{tab:notation} collects the symbols.

\begin{table}[t]
\centering
\small
\begin{tabular}{@{}p{3.4cm}p{9.6cm}l@{}}
\toprule
symbol & meaning & where \\
\midrule
$n$ & roster size; also the density window (slots) and trust depth
  (blocks) & \S\ref{sec:protocol} \\
$\quorum = \lceil 2n/3 \rceil$ & blocks needed per window & \S\ref{sec:protocol} \\
$\fmax = \lfloor (n{-}1)/3 \rfloor$ & fault budget per window & Asm.~\ref{ass:budget} \\
$\tau$ & slot length (seconds) & \S\ref{sec:impl} \\
$b$, $\Delta$ & recency bound: accept a tip at most $b$ slots old ($\Delta$ in Lean) & \S\ref{sec:verifier-rule} \\
$\Delta_{\mathrm{conf}}$ & mode 1's confirmation lag (a different parameter) & \S\ref{sec:rotation} \\
$\delta$, $\varphi = b + \delta$ & verifier clock error; tip age bound in real time & \S\ref{sec:verifier-rule} \\
$\ell$, $\sigma$ & how far back the budget looks; how far honest clocks may run ahead & Asm.~\ref{ass:budget}--\ref{ass:clock} \\
$\mathit{keyFor}(i,j)$, \code{keyIndex} & validator $i$'s public key at
  version $j$; the version a block uses & \S\ref{sec:protocol} \\
$\rho$, $T$ & rented slots and stolen-key validators per window
  (mode 3) & \S\ref{sec:results} \\
$\mathit{gen}(s)$ & scheduled key version at slot $s$ (mode 2) & \S\ref{sec:results} \\
\bottomrule
\end{tabular}
\caption{Notation.}
\label{tab:notation}
\end{table}

\subsection{Blocks and chains}

A block has six fields:
\[
B = (\mathit{slot},\ \mathit{height},\ \mathit{prev},\ \mathit{id},\
\mathit{contentsHash},\ \mathit{keyIndex}).
\]
\begin{itemize}
\item \code{contentsHash} commits to the payload (the transactions),
  which travels separately. No consensus rule looks inside it.
\item \code{keyIndex} is the version of the key that signed the block.
  Key rotation is therefore part of the chain itself.
\item \code{prev} is the parent's id, and \code{id} is the block's own
  hash.
\end{itemize}
The producer's signature travels next to the block, not inside it.

Ids follow a fixed rule, which the deployed client and the
certificate's circuit both recompute:
\[
\mathit{id} \;=\; H(\mathit{slot},\ \mathit{height},\ \mathit{prev},\
\mathit{parentSig},\ \mathit{contentsHash},\ \mathit{keyIndex}).
\]
The model's rules treat ids as opaque, but two inputs of the hash
matter later:
\begin{itemize}
\item The \emph{parent's signature} is hashed in. So anyone building a
  fake chain must produce its signatures one after another, never all at
  once. Section~\ref{sec:timed} turns this into a bound in real time.
\item The \emph{key version} is hashed in. Otherwise two blocks that
  differ only in key version would have the same id as soon as anyone
  rotated a key.
\end{itemize}

A \emph{chain} is a sequence of blocks starting at a genesis block
(height zero, no parent). Heights go up by one, slots strictly increase,
and each block points to its parent's id (\code{linksOk},
\code{genesisOk}).

\subsection{The density rule}

There are no votes, and chains are not compared by length. Instead, a
chain must be \emph{dense}: every completed window of $n$ consecutive
slots must contain at least $\quorum$ of its blocks (\code{denseSoFar};
\code{validChain} combines this with the structure checks). A window is
complete once the chain's newest block has reached its end.

In plain terms: a chain counts as evidence only if, all along its
history, at least two thirds of the roster was building it.

Density is the whole safety mechanism. Suppose two chains split and both
stay dense. In any window after the split, each chain needs $\quorum$
blocks. Since $2\quorum > n + \fmax$, some slot in that window has a
block on \emph{both} chains and belongs to an honest validator. But an
honest validator signs only one block per slot. Every safety theorem in
Section~\ref{sec:results} is built on this counting argument, with
signatures, certificates and key rotation layered on top.

One more rule uses \code{keyIndex}: along a chain, a validator's key
version never goes down (\code{keyMonoOk}). Rotation only moves forward,
and the chain's order is the agreed order of rotations. Structure,
density and this rule together are \code{validChainK}. This monotone
rule belongs to mode 1; modes 2 and 3 replace it with their own pins,
and Theorem~\ref{thm:lc} and the shipped validators use no key-version
rule at all. A deployment completes the protocol by choosing a rotation
mode (Section~\ref{sec:rotation}).

\subsection{What a node holds}

A node does not keep history. It keeps a \emph{certified chain}: a
constant-size certificate plus a tail of recent signed blocks. The
certificate vouches for everything below the tail. It exposes a
\emph{claim}: the slot, height and id of the newest block it covers,
plus a few blocks near that point so that density can be counted across
the join.

In our implementation the certificate is a \emph{recursive}
zero-knowledge proof. Each new proof checks the previous proof plus the
new blocks, so one small object covers an arbitrarily long history.

The rules use cryptography through two small interfaces that the
deployment provides:
\begin{itemize}
\item \code{SigOps}: sign and verify, plus the key map
  $\mathit{keyFor}(i,j)$, the public key of validator $i$ at version $j$
  (in the model, the map is a separate parameter, \code{KeyRegistry});
\item \code{CertOps}: verify a certificate, read its claim, and extend it
  by one checked block.
\end{itemize}

There is no key directory to trust: a verifier looks nothing up. Each
validator has one offline master key (a \emph{cold root}), created at the
deployment ceremony. A block's signature is really a bundle:
\begin{itemize}
\item the signing public key,
\item a delegation certificate, signed by the cold root, that binds this
  key to validator $i$ \emph{and version $j$}, and
\item the signature itself.
\end{itemize}
The delegation must name the version. A key that could answer for
several versions is exactly the key reuse that Section~\ref{sec:results}
must rule out.

\paragraph{What ``the signature verifies'' means.} The theorems read
signature checks the standard idealized way, as the ideal signature
functionality $\mathcal{F}_{\mathrm{SIG}}$~\cite{canettisig}: at a given
moment, a signature verifies only if someone has already produced it,
either the honest key holder or whoever holds a leaked key. A signature
scheme behaves this way exactly when it is EUF-CMA, the standard notion
of unforgeability~\cite{canettisig}. Assumption~\ref{ass:sig} states this
as an assumption (\code{hbridge}).

\subsection{The verifier rule}
\label{sec:verifier-rule}

Everything a light client does is one rule. Accept a certified chain if
and only if:
\begin{enumerate}
\item it is \emph{valid}: it passes the five checks of
  Section~\ref{sec:rules}, namely certificate, signatures, links,
  density, and key versions under the deployment's mode
  (\code{validCertifiedChain} plus the mode's key checks); and
\item it has a \emph{tail}: at least one signed block follows the
  certificate. In mode 3, and for the untimed mode-2 certificate
  theorems, the block trusted must lie $n$ deep ($2n$ under mode 3's
  per-generation budget) counted in the history the certificate attests
  followed by the tail; only the suffix-only forms need the tail itself
  to be that long; and
\item it is \emph{recent}: its newest block is at most $b$ slots old on
  the verifier's clock.
\end{enumerate}
Then trust everything $n$ blocks below the newest block.

Every certificate-level theorem in this paper is about certified chains
that pass this rule; a bare certificate, with no tail, is covered by none
of them. The model theorems take the rule as hypotheses on the tail
(\code{hLink}, \code{hLinks}, \code{hDense}, \code{hSigned}, the mode's
key checks, \code{hRecent}); the deployed Rust client checks all of it itself, the
TypeScript one all but the key checks (Section~\ref{sec:impl}).

If the verifier's clock is off by at most $\delta$ slots, the newest block
is at most $\varphi = b + \delta$ slots old in real time. The deployed
client makes the comparison; the deployment supplies only the clock
reading (Section~\ref{sec:limitations}).

Both halves are needed. Validity without recency lets in forks signed
with keys that leaked long afterwards: the long-range attack. Recency
without validity means nothing.

Choosing $b$ is a trade against the fault budget: a leaked key counts
against a window until $\varphi$ slots after it
(Assumption~\ref{ass:budget}). For example, $b = n$ with the prototype's
settings ($n = 7$, $\tau \approx 4$\,s) accepts tips up to about half a
minute old.

\section{The rules a node runs}
\label{sec:rules}

A node does three things with its certified chain: validate, produce and
select. All three are pure functions, with no I/O and no hidden state.
That is what lets one source serve as the reference implementation, the
Lean model and the circuit of the proof (Section~\ref{sec:impl}).

\paragraph{Validate} (\code{validCertifiedChain}). Accept a certified
chain only if:
\begin{enumerate}
\item the certificate verifies (\code{CertOps.verify}), one cheap check;
\item every block in the tail has a signature that verifies under the key
  of its slot's validator, at the key version the block states;
\item the tail links up: its first block extends the certificate's
  newest block (height one up, slot later, \code{prev} equal to that
  block's id), and so does each following block;
\item every window that completes inside the tail is dense, counting the
  certificate's boundary blocks plus the tail (\code{validSuffix});
\item key versions follow the mode's rule: in mode 1 a validator's
  version never goes down (\code{keyMonoOk}) and every block meets the
  pin read from the chain; mode 2 checks the schedule pin instead, and
  mode 3 the lockstep rule. This is the only check that depends on the
  mode, and Theorem~\ref{thm:lc} uses none of it.
\end{enumerate}
Check~5 is not part of \code{validCertifiedChain}, which runs checks 1
to~4 only. The key-version checks are stated for full chains
(\code{validSignedChainK$'$} and its mode 2 and 3 variants), and appear
in the certificate theorems as hypotheses on the tail. At a certificate
boundary, mode 1's check needs each validator's version floor, which the
fold certificate carries and the Rust client takes from it
(Section~\ref{sec:rotation}).

The model re-checks density from the genesis (\code{validChain}). A
running node checks only the windows that complete with each new block.
The two agree: a valid chain extended by a child of its tip is valid
exactly when every window the new block completes is dense
(\code{validChain\_snoc}).

\paragraph{Produce} (\code{produceBlock?}). Only in its own slot, a
validator builds the next block on its current tip, signs it and appends
it. It ships the result only if it passes the same \code{validSuffix}
check every other node runs. So nothing malformed that production could
emit needs to be trusted: it is rejected everywhere. The node's signing
discipline is another matter. Signing at most once per slot, a clock at
most $\sigma$ ahead, and a signing record that survives restarts are
assumed by every safety theorem (Assumption~\ref{ass:sig}(c)); a node
that breaks them is charged to the budget, like a stolen key.

\paragraph{Select} (\code{selectChain}). Switch to a candidate chain only
if it is valid and its newest block is strictly higher. Selection never
affects safety: the theorems of Section~\ref{sec:results} hold for
\emph{any} two chains a verifier accepts, however they were chosen.
Selection matters only for liveness.

\paragraph{Compaction.} From time to time, a node folds the oldest tail
blocks into the certificate, keeping the tail short, but never shorter
than the verifier rule requires. Certificates are
proofs, not authority, so anyone may run the prover; one untrusted
proving service is enough. If proving falls behind, tails simply grow
until it catches up.

The compacted chain must pass \code{validCertifiedChain} like any
other, so a bad cut is rejected, and safety needs no trust in
compaction. That compaction always finds a certificate for a valid
extension is a liveness question, and it is not proved.

\section{What a deployment must provide}
\label{sec:assumptions}

Each assumption below appears as a hypothesis of the theorems that use
it; Lean's axiom check confirms that nothing else is assumed.
Theorem~\ref{thm:lc} is stated at the real time $R$ at which the
verifier checks; the rotation theorems use the verifier's clock reading
$\mathit{now}$; Assumption~\ref{ass:clock} links the two. Mode 2 keeps
these assumptions, reading key exposure at the scheduled key versions.
Mode 3 uses adjusted forms (Section~\ref{sec:rotation}).

\paragraph{The attacker.} It controls message delivery completely;
timely delivery matters only for liveness (Assumption~\ref{ass:delivery}).
It obtains signing keys, by theft or by controlling a signer for a while,
at any time, within the budget of Assumption~\ref{ass:budget}, and a
leaked key signs anything, for any slot, at any time. Its timing is
unconstrained. Two clocks matter: the verifier's
(Assumption~\ref{ass:clock}) and the honest signers', which run at most
$\sigma$ slots ahead (Assumption~\ref{ass:sig}(c)). It cannot sign with
keys that have not leaked, find hash collisions between real blocks, or
forge certificates (Assumptions~\ref{ass:sig}--\ref{ass:cert}).

\begin{assumption}[Exposure budget]\label{ass:budget}
A key is \emph{exposed} at a given moment if someone other than its
honest holder can sign with it then: the validator is malicious, the key
was stolen, or the attacker controls the signer. Exposure can start and
stop: a stolen key stays exposed, while control of a signer may come and
go.

The budget: in any $n$ consecutive slots, at most $\fmax$ slots have a
key that is exposed at some moment in a stretch of real time around that
window, from $\ell$ slots before it to $\varphi$ slots after it.
Exposure outside that stretch does not count against the window. For
the usual $n = 3\fmax + 1 \ge 4$, a lookback $\ell = n - 1 + \sigma$ is
enough; the general condition is in the Lean. For $n \le 3$, $\fmax = 0$:
the results hold but tolerate no exposure.
\quad\textit{(Lean: \code{ExposureBounded}.)}
\end{assumption}

\begin{assumption}[Signatures]\label{ass:sig}
This is about \emph{admissible} blocks, those the verifier's own check
would accept: the signature verifies and the id is correctly formed
(Section~\ref{sec:protocol}).
\begin{enumerate}
\item[(a)] \emph{Custody.} While a key is not exposed, only its honest
  holder signs with it.
\item[(b)] \emph{Unforgeability.} A signature the verifier accepts at real
  time $R$ was produced by time $R$, by the honest holder or by whoever
  holds the leaked key (\code{hbridge}).
\item[(c)] \emph{Honest signing.} An honest validator signs at most one
  block per slot, ever, and its clock runs at most $\sigma$ slots ahead.
\end{enumerate}
\quad\textit{(Lean: \code{SigningExecutionOn}, \code{HonestClockOn},
\code{hbridge}.)}
\end{assumption}

For part (c), an
honest validator keeps a persistent record of what it has signed, apart
from the key material and outliving it: a validator that loses its key
and signs again under a new version must still refuse every slot it
already signed. A slow clock costs only liveness.

One more property is assumed, and it is why ids include the parent's
signature: a block can be signed only after its parent exists
(\code{chain\_order}, a field of \code{SigningExecutionOn}). An id rule
without the parent's signature would make this false, and no check would
notice. Junk that no verifier accepts, such as a
block with a copied id, is outside the assumption, so it cannot break it.

\begin{assumption}[Hash collision resistance]\label{ass:hash}
No two different admissible blocks that actually occur (or the genesis)
have the same id.
\quad\textit{(Lean: \code{SigningExecutionOn.id\_inj},
\code{SignedHashInjective} over \code{SignedFormed}.)}
\end{assumption}

Restricting the statement to blocks that actually occur turns ``no one
finds a collision'' into a fact about ids. Mode 3 and the untimed
mode 2 theorems use it over blocks that verify under the version they
state and have correctly formed ids (\code{SignedFormed}). The second
half is needed: a thief holding one key can sign two blocks with the
same copied id. \code{IdInjective} over a record of the blocks that
occur (\code{ChainInRecord}) is the form \code{same\_block\_same\_prefix}
uses to extend agreement on one block to its whole prefix.

\begin{assumption}[Certificate soundness]\label{ass:cert}
A certificate that verifies vouches for a claim built from the genesis
by adding one block at a time. Each step checks the links, the density
of newly completed windows, and that the new block is admissible
(Assumption~\ref{ass:sig}). Modes 2 and 3 add their key-version check
to each step (Section~\ref{sec:rotation}).
\quad\textit{(Lean: \code{GroundedCert}, \code{GroundedCertSched},
\code{GroundedCertLock}.)}
\end{assumption}

The prototype's fold certificate (Section~\ref{sec:impl}) runs one
verified step per block, and the step is proved to build exactly these
claims (\code{cert\_fold\_k\_sound}, \code{cert\_fold\_sched\_sound},
\code{cert\_fold\_lock\_sound}). For the deployed Rust client the two
are composed (Section~\ref{sec:impl}): its theorems hold for every
history that folds to the certificate's public state, and knowledge
soundness supplies one (\code{rust\_deployed\_ks}). For that client the
assumption becomes the circuit's knowledge soundness, that its gadgets
compute what the verified code computes, and that it routes each
witness as the fold relation says. The model theorems and the
TypeScript client take it as stated.

The assumption does not posit a valid chain behind the certificate: a
proved lemma rebuilds it from the steps (\code{groundedCert\_history}
and its mode-2 and mode-3 twins), and the certificate theorems conclude
about that history followed by the tail. Theorem~\ref{thm:lc}'s
certificate needs no signed genesis; the full-chain validators of every
mode, and the certificate derivations and fold steps of modes 1--3,
check the genesis like any other block, signed by validator~0. Every
certificate model places the genesis at slot $0$.

\begin{assumption}[Clock]\label{ass:clock}
The verifier's clock reads $\mathit{now}$, and real time is at most
$\delta$ slots ahead of it; $\delta$ covers clock error and network delay.
\end{assumption}

The verifier accepts a newest block when $\mathit{now} \le
\mathit{tip.slot} + b$; Theorem~\ref{thm:lc} then applies with
$\varphi = b + \delta$. A larger $b$ tolerates slower delivery, at the
price of a longer stretch in which exposure counts against each window.

\begin{assumption}[Honest delivery, for liveness only]\label{ass:delivery}
The honest validators of each newly completing window were online, and
their blocks reached the next producer in time to be built on.
\quad\textit{(Lean: \code{HonestBlocksCover}.)} The Lean form also
assumes that every honest block is on the one chain the theorem is
about: a fork-choice property, which we do not model
(Section~\ref{sec:liveness}).
\end{assumption}

\section{What the theorems guarantee}
\label{sec:results}

This section states the results in plain words; each name in brackets
points to the exact statement. We cover light-client safety
(Theorem~\ref{thm:lc}), the fact it rests on, that no valid chain runs
ahead of real time (Theorem~\ref{thm:forge}), the rotation modes, and
liveness.

\paragraph{Two kinds of attacker.} The \emph{exposure} attacker of
Section~\ref{sec:assumptions} obtains keys at any real time, within the
budget. Theorems~\ref{thm:lc} and~\ref{thm:forge} face it, counting a
slot as exposed if any of its validator's keys is. Mode 2's theft
theorem faces the same attacker, but reads exposure at the scheduled key
version. The \emph{key-stealing} attacker of Section~\ref{sec:rotation}
is counted slot by slot: rented slots plus stolen key versions that are
still live. Mode 3 and the untimed mode 2 theorems face it.

\paragraph{How far each result reaches.} Every result is a Lean theorem
about the model. Theorem~\ref{thm:lc} and its mode-1, mode-2 and mode-3
forms are also proved for the deployed Rust client, from the certificate
it verifies to the blocks it accepts
(\code{rust\_deployed\_agreement}, \code{rust\_deployed\_k\_agreement},
\code{rust\_deployed\_sched\_agreement},
\code{rust\_deployed\_lock\_agreement}; Section~\ref{sec:impl}), and
Theorem~\ref{thm:lc} for the TypeScript client
(\code{ts\_deployed\_agreement}). Those clients recompute block ids and
check recency themselves. The other results are proved for the model
only.

\subsection{Light-client safety}
\label{sec:lc}

\begin{theorem}[Recent chains share their deep blocks; Lean
\code{timed\_light\_client\_safety}]\label{thm:lc}
Fix $n \ge 1$, a genesis $G$, a freshness $\varphi$, an honest clock
lead $\sigma$, and a lookback $\ell$ with $\ell \ge n - 1$ and
$\ell \ge n + \fmax + \sigma + 1 - \quorum$. Assume:
\begin{itemize}
\item \emph{Signing} (\code{SigningExecutionOn}): Assumption~\ref{ass:sig}
  (a) and (c) and Assumption~\ref{ass:hash}, for admissible blocks:
  while a key is unexposed it signs at most one admissible block per slot
  (\code{honest\_once}); a block is signed only after its parent exists
  (\code{chain\_order}); no two admissible blocks share an id
  (\code{id\_inj});
\item \emph{Clocks} (\code{HonestClockOn}): honest clocks run at most
  $\sigma$ slots ahead;
\item \emph{Budget} (\code{ExposureBounded}~$n$~$\ell$~$\varphi$):
  Assumption~\ref{ass:budget};
\item \emph{Unforgeability} (\code{hbridge}): Assumption~\ref{ass:sig}(b),
  the verifier checking at real time $R$;
\item \emph{two accepted chains}, each a certificate (\code{GroundedCert})
  plus a non-empty tail the verifier accepts (Section~\ref{sec:rules}),
  both recent: $R \le \mathit{tip.slot} + \varphi$.
\end{itemize}
Then the two full chains (\code{GroundedHistory}) agree on every block at
least $n$ blocks below the lower tip.
\end{theorem}

It also covers a chain the node \emph{produced}, since production ships
nothing that would fail the validator. The theorem compares two chains
checked at one real time $R$. Clients that check at different times are
covered when the earlier chain is still recent at the later time: apply
it there to both. Anything stronger needs the honest chain to keep
growing between the two checks, which is network liveness
(Section~\ref{sec:liveness}); we state no such theorem.

\paragraph{Why it holds.} The classic quorum count, with the timing made
explicit. In the last full window below the lower tip, each chain has at
least $\quorum$ blocks, so at least $2\quorum - n \ge \fmax + 1$ slots
have a block on both. If all those pairs differ, each slot's key was
exposed at one of the two signing times (\code{honest\_once}). Each
signature was made by time $R$, at most $\varphi$ slots after the window,
and by Theorem~\ref{thm:forge} no earlier than $\ell$ slots before it.
That is $\fmax + 1$ exposed slots within one budget stretch, which the
budget forbids. So the chains share a block in the window, and shared ids
and collision resistance make them agree below it.

\paragraph{Why the lookback.} Without one, every exposure up to $\varphi$
slots after a window counts against it, however long ago it happened.
Under that cumulative budget (\code{ExposureBoundedEver}) agreement holds
with no assumption on clocks, for full chains whose non-genesis blocks
are admissible and exist by time $R$
(\code{exposure\_agreement\_ever\_on}), and for a certificate plus a
tail (\code{exposure\_certified\_agreement\_ever\_on}). The lookback lets a
cleaned-up signer stop counting; the honest clock bound is its price.

\subsection{No chain runs ahead of real time}
\label{sec:timed}

The recency check needs one more fact: an attacker must not be able to
build a chain early whose tip will look recent later. It rests on
\code{chain\_order}. A block's id includes its parent's signature, so a
block can be signed only after its parent exists; predicting someone
else's signature without their key would be a forgery.

\begin{theorem}[No valid chain runs ahead of real time; Lean
\code{no\_early\_signing}]\label{thm:forge}
Assume the signing, clock and budget assumptions of
Theorem~\ref{thm:lc}, including the lookback condition. Take a valid
chain from $G$ whose blocks are admissible and all exist by some real
time. If one of its blocks $B \ne G$ was signed at real slot $r$, then
$B.\mathit{slot} \le r + \ell$. With $\ell = n - 1$ (for
$n = 3\fmax + 1 \ge 4$ and $\sigma \le \fmax - 1$), no valid chain gets
more than $n - 1$ slots ahead of real time.
\end{theorem}

\paragraph{Why it holds.} By induction along the chain. If $B$ was signed
more than $\ell$ slots early, its signer was not honest (honest clocks
lead by at most $\sigma \le \ell$), so its key was exposed at $r$. The
window ending at $B$ holds at least $\quorum$ blocks, all signed by time
$r$. An honest one has slot at most $r + \sigma$; an exposed one was, by
induction, not signed early either, so its exposure counts against the
window. Either way, too few blocks remain to reach $\quorum$.

\subsection{Key rotation}
\label{sec:rotation}

Each block states its key version (\code{keyIndex}). A deployment adds
one rule, the \emph{pin}: a lower bound on the version a block may state
(exact per grid window in mode 3). The pin is what makes a retired key
\emph{dead} rather than merely old. The modes differ only in what sets
it: the chain itself (mode 1), the slot number through a fixed schedule
(mode 2), or a roster-wide counter read off the signatures (mode 3).
That choice decides what a light client must hold.

\paragraph{Lost keys, in every mode.} A destroyed key is held by no one,
so losing it exposes nothing and charges nothing to the budget, provided
the signing record survives (Assumption~\ref{ass:sig}(c)). This is how a
deployment reads Assumption~\ref{ass:budget}, not a separate theorem:
Theorem~\ref{thm:lc}'s assumptions are unchanged, and restricting it to
each mode's validator gives \code{keyrot\_loss\_agreement} (mode 1,
Theorem~\ref{thm:keyloss}), \code{sched\_loss\_agreement} and
\code{lockstep\_loss\_agreement}. In mode 3's slot-by-slot count, with
nothing stolen only rented slots count (\code{badSched\_lossOnly}).

\paragraph{Stolen keys.} A stolen key of \emph{any} version signs
anything, for any slot, forever: we assume no forward
security~\cite{forwardsec}. Mode 2 needs nothing beyond
Theorem~\ref{thm:lc}'s assumptions, read at the schedule. Mode 3, and an
older untimed form of mode 2, count per window of $n$ slots at most
$\rho$ \emph{rented} slots plus at most $T$ validators with a \emph{live}
key stolen, $\rho + T \le \fmax$. A rented slot is controlled for that
slot only; a rented device made to sign for other slots counts as
stolen. A stolen key is live if its version is at least the one in
force, so it stops counting once its retirement takes force
(\code{badSched}). Looking back, the count can still charge windows
before the theft. It also needs \emph{version independence}: stealing
version $j$ reveals nothing about a later $j'$, which key-insulated
signatures~\cite{keyinsulated,keyinsulatedsig} provide.

\begin{assumption}[Honest custody, per key]\label{ass:rotation-honest}
Used by mode 3 and the untimed mode 2. For a slot $s$ that is not rented
and a never-stolen version $j$ of its validator $i$: every block for $s$
that verifies under $j$ was signed by $i$'s own device; that device signs
at most one block per slot, ever; and in mode 3 it states the version the
mode assigns to the slot.
\quad\textit{(Lean: \code{SchedUnforgeable}, \code{SchedCoreUnforgeable};
the third item is the \code{declared} field of mode 3's packages.)}
\end{assumption}

It is assumed rather than derived because it is stated per key version,
and one slot can have several eligible versions, only one of which
verifies; Theorem~\ref{thm:lc}'s exposure model counts producers, not
versions (\code{badSched\_single\_key\_safe\_not\_enough} is a sanity
check exhibiting such an instance). The untimed theorems have no clock
model: they read the recency check at one $\mathit{now}$, shared by both
clients. Their two-reading forms let each client check at its own
$\mathit{now}$, with the budget read at the earlier one
(\code{sched\_recent\_certified\_history\_agreement\_horizon\_two\_readings},
\code{lockstepGen\_recent\_certified\_history\_agreement\_two\_readings}).

\subsubsection*{Mode 1: rotate when you want}

A validator rotates whenever it chooses, by signing under a higher
version. Its version \emph{in force} (\code{inForce}) is the highest it
has used in the chain's \emph{confirmed} part, blocks at least
$\Delta_{\mathrm{conf}}$ slots old, and no block may state less
(\code{validSignedChainK$'$}). On a full chain this adds nothing to the
monotone rule, for any $\Delta_{\mathrm{conf}}$. At a certificate
boundary the verifier has no history, so the certificate must supply
each validator's floor. The fold certificate does: each step checks the
block's version against its producer's floor, raises it, and exposes all
$n$ floors. The deployed Rust client decodes them from the
certificate's public inputs and checks the tail against them, and both
full chains then satisfy the monotone rule
(\code{rust\_deployed\_k\_agreement}), so a certificate client does
enforce that a retired key is dead. With $\Delta_{\mathrm{conf}} = n$, a rotation
takes effect in under half a minute at the prototype's settings.

A validator that loses its key signs its next block under a higher
version, whose key the deployment must already have delegated from the
cold root (\code{liveness\_produce\_blockK}).

\begin{theorem}[Mode 1 under key loss; Lean
\code{keyrot\_loss\_agreement}]\label{thm:keyloss}
Under the signing, clock, budget and unforgeability assumptions of
Theorem~\ref{thm:lc}, with a block admissible
when it verifies under the version it states (\code{SignedDeclared}) and
its id is correctly formed: two chains from $G$ that pass the mode-1
validator (\code{validSignedChainK$'$}, any confirmation depth), with
correctly formed ids (\code{hFormed}) and recent tips, agree on every
block at least $n$ below the lower tip.
\end{theorem}

\paragraph{Not safe against theft.} The pin keeps a retired key dead
only on chains that contain its retirement. Validators stop guarding
retired keys; an attacker who later steals those of a quorum builds a
dense fork from a point before the rotations, with recent slots. It
passes the mode-1 validator and the recency check, and no budget on the
real chain rules it out. A deployment that must survive theft runs mode
2 or 3.

\subsubsection*{Mode 2: rotate on a schedule}

The version in force at slot $s$, the \emph{generation}, is fixed in
advance: $\mathit{gen}(s) = \lfloor s/E \rfloor$ for an era length $E$.
It is the same on every chain and built into the verifier, never taken
from the network. The validator adds one check,
$\mathit{gen}(\mathit{slot}) \le \mathit{keyIndex}$
(\code{validSignedChainSched}). Since the pin reads nothing from any
chain, a fork cannot disagree about which keys are retired.

Slot $s$, of validator $i = s \bmod n$, is \emph{exposed} at real time
$r$ (\code{schedExposed}) when the attacker controls $i$
($\mathit{Controlled}(i, r)$), or someone other than $i$ can sign with a
key of $i$ of version $j \ge \mathit{gen}(s)$ ($\mathit{Stolen}(i, j, r)$).
Both may change over time arbitrarily. A block is \emph{admissible} when
it verifies under the version it states, that version is at least
$\mathit{gen}$ of its slot (\code{SchedAdmissible}), and its id is
correctly formed; a block signed with a key already retired for its slot
is not.

\begin{theorem}[Mode 2 under key theft; Lean
\code{sched\_exposure\_agreement},\hfil\penalty0
\code{sched\_exposure\_certified\_agreement}]\label{thm:sched}
Fix $\mathit{gen}$ and any $\mathit{Controlled}$ and $\mathit{Stolen}$.
Under the signing, clock, budget and unforgeability assumptions of
Theorem~\ref{thm:lc}, with exposure and
admissibility read as above, two chains from $G$ that pass the mode-2
validator, with correctly formed ids and recent tips, agree on every
block at least $n$ below the lower tip. The same holds for two
certificate-plus-tail presentations whose certificates check the pin at
each step (\code{GroundedCertSched}, \code{GroundedHistorySched}).
\end{theorem}

Theorems~\ref{thm:lc} and~\ref{thm:forge} are stated for an arbitrary
per-block test, and Assumption~\ref{ass:sig} is needed only of blocks
that pass it; mode 2 picks admissibility. Required of every signed
block, the assumption would fail as soon as a stolen key signed a block
with a made-up parent.

\paragraph{What the budget charges.} A stolen version $j$ exposes only
slots of generation at most $j$. It counts against every window whose
stretch reaches the theft, until era $j$ ends: a theft up to
$n + \varphi - 1$ slots after era $j$'s last slot still counts against
its last windows. A key retired more than $n + \varphi - 1$ slots ago exposes nothing, stolen or
not, with no erasure. A key made early and stolen early counts until its
era ends; a stolen cold root exposes its validator in every later
window, like a malicious one. So careful provisioning is not an extra
assumption; it is how an operator stays within the budget.
$\mathit{Stolen}$ must list every version the attacker can sign with;
version independence keeps that list as short as what was taken.

\paragraph{Operator duties.} Keep one offline cold root per validator,
and delegate each generation's key just in time for its era; a theft then
costs only the rest of that era. Longer eras mean fewer ceremonies but
longer exposure of a stolen live key. A validator that loses its live key
is offline until its next generation, its slots counting against $\fmax$.
The model has one key per version (\code{KeyRegistry}); a deployment
that re-delegates a fresh key for the same generation must count both in
$\mathit{Stolen}$. What breaks first: a schedule taken from the network
(a zero schedule would readmit every retired key); theft of the current
generation, which counts until its era ends; a stolen cold root, which
costs its validator for good.

\paragraph{The untimed form.} An older form has no real-time model of
signing. Under
Assumption~\ref{ass:rotation-honest} and the $\rho + T$ count, which
charges a stolen live key from the start of time, two accepted chains
agree on the block $n$ below each tip
(\code{sched\_recent\_tip\_ancestor\_agreement\_horizon}, and
\code{sched\_recent\_certified\_history\_agreement\_horizon} over a
certificate). Theorem~\ref{thm:sched} supersedes it.

\paragraph{A mode-2 deployment.} At setup each validator makes its cold
root offline and the deployment publishes $\mathit{gen}$, with $E$ a
week of slots, say. Every week each cold root signs one delegation for
the coming generation and goes back in the safe. Wallets hold the
genesis, a clock and $\mathit{gen}$, and may sleep for months. What must
hold: Assumption~\ref{ass:sig} for admissible blocks, collision
resistance, the clock, and the budget over mode-2 exposure, sized for a
stolen live key counting the rest of its week plus the window's
stretch. What you get: by Theorem~\ref{thm:sched} and the parent-id
links (\code{same\_block\_same\_prefix}), any two accepted recent chains
agree on everything at least $n$ below the lower tip, whatever happened
while the wallet slept. The prototype's fold certificate checks the pin
at each step, the deployed Rust client checks it on the tail, and
Theorem~\ref{thm:sched} holds for that client end to end
(\code{rust\_deployed\_sched\_agreement}).

\subsubsection*{Mode 3: rotate together, at any time}

There is no schedule. The roster moves one shared generation counter
forward whenever it decides to, after an incident or by governance.
Validity enforces only the discipline of the switch (\code{noMixing},
\code{validSignedChainLock}): within each grid window (slots
$[Wn,\,Wn{+}n)$, a fixed grid), all blocks state the same generation,
and it never goes down from one grid window to the next. The verifier
reads each window's generation off the signatures.

\begin{theorem}[Lockstep safety, per generation; Lean
\code{lockstep\_client\_safety\_gen},\hfil\penalty0
\code{lockstepGen\_recent\_tip\_ancestor\_mem}]\label{thm:lock-gen}
Assume the lockstep validator; honest signatures state their grid
window's counter; hash collision resistance;
Assumption~\ref{ass:rotation-honest}; at most $\rho$ rented slots per
window and, for each generation $j$ separately, at most $T$ validators
whose generation-$j$ key is stolen, $\rho + T \le \fmax$; both tips
recent, $\mathit{now} \le \mathit{tip.slot} + \Delta$; and chains longer
than $2n$ blocks (\code{hLong}). Then two accepted chains agree on the
block $2n$ below each tip at equal heights (\code{hTipHeight}); at
unequal heights, the lower chain's block $2n$ deep is a block of the
taller one, at least $2n$ deep there. Through the parent-id links they
agree on their whole common prefix
(\code{lockstepGen\_recent\_genesis\_agreement}).
\end{theorem}

\paragraph{Why it holds.} A pinning theorem
(\code{lockstep\_window\_declares\_rosterGen}) shows that the last
complete grid window below the lower tip states one generation on both
chains, so the quorum count runs at a single generation
(\code{window\_shared\_prefix}). The tip's own window is incomplete,
hence depth $2n$, exactly $n + ((t{+}1) \bmod n)$ for lower tip slot $t$
(\code{lockstepGen\_shared\_prefix\_sharp}).

\begin{theorem}[Lockstep safety, lifetime budget; Lean
\code{lockstep\_client\_safety}]\label{thm:lock}
With the $\rho + T$ budget read at the counter one grid window back,
instead of per generation, the conclusions of Theorem~\ref{thm:lock-gen}
on the block below each tip, at equal and at unequal heights, hold at
depth $n$ (\code{lockstep\_recent\_tip\_ancestor\_mem} at unequal
heights), for chains longer than $n$ blocks.
\end{theorem}

Here every complete grid window of an accepted recent chain states
exactly the roster's counter (\code{lockstep\_declares\_rosterGen}), so
the untimed mode-2 results carry over. But theft is not time-stamped and
the counter starts at its genesis value, so window~$0$ already counts
every validator that \emph{ever} has a key of a used generation stolen:
$T$ caps key losses over the deployment's life, two at $n = 7$.
Rotation does not reset it. The first budget gives the second
(\code{LockstepPackage.toGen\_of\_low}), so the per-generation one is
the weaker requirement and the one to plan with.

\paragraph{Operator duties.} Two, and they are this mode's real price.
All honest signers follow the one counter; a split among them is an
operational failure the theorems do not cover. And retired keys are
really erased at each switch. Erasure earns its keep formally: with
theft time-stamped, it says no generation is stolen after the roster has
moved past it (\code{ErasureTimedLock}). That turns the count a
deployment attests while a generation is live into the per-generation
count (\code{genBound\_of\_preRetirementBound}; end to end,
\code{lockstep\_client\_safety\_timed}; \code{LockstepPackageTimed.toGen}
gives every form of Theorem~\ref{thm:lock-gen}). So a deployment budgets
$T$ validators whose generation-$j$ key is stolen \emph{before $j$
retires}. Theft times are read on the grid, so ``erased at the switch''
means before the grid window in which the counter moves starts. Keys in a hardware module, or only ever in
RAM, qualify; flash storage, which does not erase by overwriting, does
not.

\paragraph{A mode-3 deployment.} Each validator signs from a hardware
module that destroys retired generations on a switch; the genesis states
the counter's starting value. Switches are agreed out of band, and every
honest signer moves at the agreed grid-window boundary. Clients hold the
genesis and a clock, nothing else, may sleep for any time, and sync from
a certificate plus a tail. The certificate carries the claim and the tip
generation; for $n \ge 2$ the claim's tail determines the generation
(\code{groundedCertLock\_gen\_unique}). What you
get: any two accepted chains agree on the lower chain's confirmed part,
with no schedule and no anchor. For certified chains,
\code{lockstep\_recent\_certified\_suffix\_agreement} and
\code{lockstepGen\_recent\_certified\_suffix\_agreement} cover tail
blocks at least $n$ (or $2n$) tail blocks below each tail's tip, and
their \code{\_history\_} forms the rebuilt history plus a non-empty tail.
Under the lifetime budget the
certificate's counter is pinned to the roster's once its window
completes (\code{lockstep\_cert\_gen\_pinned}). These need a certificate
that carries the generation and checks the lockstep rule
(\code{GroundedCertLock}); the prototype's fold certificate does, the
deployed Rust client checks the lockstep rule on the tail, and
Theorem~\ref{thm:lock-gen} holds for that client end to end, at depth
$2n$ in the full chains (\code{rust\_deployed\_lock\_agreement}).

\paragraph{The limits of this subsection.} Mode 1 has no guarantee under
theft. The TypeScript client runs no mode's key rule; the Rust client
runs all three. Mode 3 and the untimed mode 2 rest on
Assumption~\ref{ass:rotation-honest}; mode 2's theft theorem does not.

\subsection{Liveness, and the slot time}
\label{sec:liveness}

\begin{theorem}[Liveness; Lean \code{production\_liveness},
\code{global\_liveness}]\label{thm:live}
Assume at most $\fmax$ slots in any $n$ belong to failing validators
(\code{ByzantineBounded}), and honest delivery
(Assumption~\ref{ass:delivery}). Locally, an honest validator holding an
accepted chain produces a block in its slot, and the extended chain
passes the same validator everyone runs. Globally, from a well-formed
genesis, a chain with one block in each slot of an increasing list
passes the validator if at most $\fmax$ slots of every window are missing
from the list.
\end{theorem}

Each mode's validator also lets an honest producer through: mode 1
unsigned (\code{liveness\_produce\_blockK}); mode 2 signed, if the
signature verifies at the stated version and that version meets the schedule and the producer's floor
(\code{sched\_liveness\_produce\_block}); mode 3 signed, if it is at
least every version on the chain and equals those already in the
producer's grid window (\code{lockstep\_liveness\_produce\_block}). An attacker
can still slow a mode down through these conditions (a stolen key used
at a higher version raises its validator's floor), but not break safety.

The global form is a density fact about one chain in the unsigned
model; reading it as liveness of a synchronous run is ours, since the
theorem models no network, delivery or adversary.

\paragraph{The slot time.} The slot length $\tau$ is set by prover cost,
not network latency. A slot must be long enough for the proving pipeline
to keep its uncovered tail at most $n$ blocks forever. Lean proves a
threshold that is sufficient, from a backlog of at most $n/2$ blocks, and
necessary over the pipeline's steady states when proving one block takes
less than a slot (\code{slot\_time\_sufficient},
\code{slot\_time\_necessary}).

\section{Implementation and measurements}
\label{sec:impl}

\paragraph{One source, three uses.} The reference implementation is one
Rust file, \code{rust/src/lib.rs}, in the subset of Rust that Charon +
Aeneas translates to Lean~4.

\emph{The native validator.} Its translation is proved sound against the
model: if the code accepts a chain, the chain is valid in the model
(\code{rust\_valid\_chain\_sound} and companions; mode~1:
\code{rust\_valid\_chain\_k\_sound}).

\emph{The deployed client.} \code{validate\_deployed} is what a light
client runs. It verifies the certificate and decodes its claim from the
certificate's public state. It checks the tail's signatures, its links
and the density of the windows it completes. It recomputes every tail
block's id, with each parent's signature threaded in, starting from the
certificate's tip signature. And it checks that the newest block is at
most $\Delta$ slots old at the client's clock reading. Its per-mode forms
add the mode's key rule on the tail: the floors decoded from the
certificate (mode~1), the schedule pin (mode~2), or the lockstep rule
(mode~3). The theorems of Section~\ref{sec:results} hold for it with
nothing assumed about the claim, the ids, the recency or the tail's key
rule. The clock hypothesis is $R \le \mathit{now} + \delta$ with
$\Delta + \delta \le \varphi$
(\code{rust\_deployed\_agreement} and its three twins). Its producer,
\code{produce\_block\_deployed}, extends a tail the client accepts to
one it accepts again, given the slot is its own and after the tip, and
every window the new block completes is dense (\code{rust\_produce\_block\_deployed\_live}). This needs size
bounds: slots, the roster, and the decoded tail plus the suffix below
$2^{62}$, and a decoded tail no newer than its tip. Each excludes a
64-bit overflow, which the checked arithmetic turns into a rejection.
\code{FoldCrypto} (\code{rust-keyrot/src/fold\_crypto.rs}) instantiates
the client's cryptography with the fold certificate and XMSS, and a
test runs the client end to end on a folded chain.

\emph{The circuits.} The same source has a validator generic over its
backend (\code{valid\_chain\_be}); over plonky2 it checks a whole chain
at once. Its translation is proved sound for the links-and-density rule,
for every $n < 2^{63} - 1$ (\code{valid\_chain\_be\_validChain}), and,
conversely, it accepts every chain the model accepts with slots and
heights below $2^{62}$ (\code{valid\_chain\_be\_complete}).

The \emph{fold certificate} (\code{rust-keyrot/src/fold\_circuit.rs}) is
the one a client relies on. Each proof verifies the previous proof of the
same circuit, or starts from the genesis, a circuit constant, and runs
one step from the verified crate on one block (\code{cert\_step\_k\_be}
and its twins): links, the windows the block completes (against a
buffer of the last $n-1$ blocks' slots), and the mode's key rule. It
computes the block's id itself, as a Poseidon hash of the block's fields
and the previous state's tip signature, with every 64-bit value split
into canonical 32-bit halves. It verifies the block's signature
subproof. The subproof's message must be that id. Its key must be the
registry's key for the producer at the version the id hashes, under a
committed registry root. The digest of the signature it verified becomes
the new tip signature. The public inputs are the fold state: tip, tip
signature, slot buffer, and floors (mode 1) or generation (mode 3). The
steps are proved to build exactly the model's certificate claims
(\code{cert\_fold\_k\_sound} and twins), and to accept every chain the
mode's full-chain rule accepts, within the same size bounds
(\code{cert\_fold\_k\_complete}, \code{cert\_fold\_sched\_complete},
\code{cert\_fold\_lock\_complete}). Tests check the circuit's state
against the native fold at every block. They also check that a rejected
block, a forged previous state, a wrong id, parent signature or key
version, and a signature by the wrong key or over the wrong block have
no proof. Gadget faithfulness (including the step's witnessed division),
the circuit's routing, and the proof system's knowledge soundness are
the remaining trust assumptions of the fold certificate.

\emph{A TypeScript twin.} \code{moltPetit.ts} reaches Lean through an
independent compiler, Thales, with its own soundness bridge
(\code{ts\_timed\_certified\_agreement}). Its deployed form
\code{validateDeployed} recomputes ids and checks recency like the Rust
client, but checks no key rule (\code{ts\_deployed\_agreement}). The vendored output corrects
recorded compiler bugs by hand (\code{tools/thales-reemission/}); those
corrections are part of the trusted base. One, a match binder that
shadowed the producer's \code{contentsHash} parameter, was found only by
outside review; no theorem covered that function. The TypeScript
\code{verify} receives slot, height, parent and id, not the whole block,
so key version and contents hash are bound only through the id.

\paragraph{What is trusted.} The Lean kernel; the pinned translators
(Charon + Aeneas, Thales with its hand corrections) and the circuit
gadgets, with the Rust crate built with overflow checks on, as its
release profile sets; the node's signing loop (the once-per-slot record,
the clock), whose guarantees every safety theorem assumes; and that the named assumptions match the deployed primitives.
Cryptography enters only as theorem assumptions.

\paragraph{Measurements.} The prototype (\code{rust-keyrot/}) pairs the
consensus rule with a minimal application, a committed key set that
changes by signed rotations. Proof system: plonky2~\cite{plonky2} over
Goldilocks with Poseidon~\cite{poseidon}; signatures:
SPHINCS+~\cite{sphincsplus,fips205} over Poseidon (\code{128s});
hardware: one Intel i5-13500 (14 cores), wall-clock. Table~\ref{tab:bench}
comes from the whole-chain circuits, which measure cost only: their
signature subproofs are not tied to blocks or keys. ``Circuit rows''
measures what the prover proves; the \emph{recursion floor} is the fixed
cost of checking one proof inside another.

\begin{table}[t]
\centering
\small
\begin{tabular}{lrr}
\toprule
operation & time & circuit rows \\
\midrule
consensus validator alone (8-block chain) & $0.035$\,s & $2^{9}$ \\
one in-circuit signature verification (subproof) & $0.36$\,s & $2^{13}$ \\
verifying one subproof (recursion floor) & $0.19$\,s & $2^{12}$ \\
certificate, 8 signatures inline & $6.6$\,s & $2^{16}$ \\
certificate, 8 signatures as recursive subproofs & $3.0$\,s & $2^{15}$ \\
2048 subproofs, balanced tree (critical path) & $5.0$\,s & $2^{13}$/node \\
\bottomrule
\end{tabular}
\caption{Prototype measurements. The verified consensus rule is under
$2\%$ of any certificate that carries signatures (row 1 against rows
4--6). Signatures dominate, and they recurse well: a balanced tree keeps
latency logarithmic in the signatures per slot ($5.0$\,s for $2048$,
against about $2048 \times 0.36 \approx 740$\,s one after another).}
\label{tab:bench}
\end{table}

The fold certificate (\code{BENCH\_FOLD=1}, $n = 7$, all modes) was
measured on a different, shared machine (AMD Ryzen~5 5600X, loaded), so
we report it against the recursion floor of the same run. A step without
the signature is $2^{13}$ rows and about $2\times$ the floor; with the
SPHINCS+ subproof and registry path, $2^{14}$ rows and $2$--$4\times$.
A fold proof verifies in about 5\,ms and is 130\,KiB at any chain
length. Scaled by Table~\ref{tab:bench}'s floor, a step would cost about
$0.4$--$0.8$\,s on the reference machine plus the $0.36$\,s subproof, an
estimate, not a measurement.

So consensus is nearly free next to the signatures, and the slot time
follows: $\tau \ge \mathit{perBlock} + 2\,\mathit{baseline}/n$
(\code{recommendedSlot}). With the per-block cost one subproof plus its
check ($0.55$\,s) and, conservatively, the whole certificate as the
fixed part, $n = 7$ gives $\tau \ge 1.4$\,s with recursive signatures and
$\tau \ge 2.4$\,s inline. We use $\tau = 4$\,s. Where stateful signatures
are acceptable, XMSS~\cite{xmssrfc} over the same gadgets shrinks the
in-circuit verifier from $2^{13}$ to $2^{10}$ rows and makes certificates
about $10\times$ faster, with one-time keys whose bookkeeping the
consensus layer already does.

\section{Related work}
\label{sec:related}

\paragraph{Density in proof-of-stake.} Ouroboros Genesis~\cite{genesis}
resolves deep forks by comparing how fast each grew after their last
common block, which means replaying history. Samasika~\cite{samasika},
Mina's consensus, carries a constant-size summary of minimum window
density, the closest earlier pairing of density with succinct proofs.
In every Ouroboros variant density decides \emph{fork choice}, so
finality is probabilistic; in Molt Petit it is a \emph{validity rule},
so finality is certain at depth $n$. The Ouroboros line is also
stake-weighted with private VRF-based leaders from
Praos~\cite{praos,ouroboros} on, and argued by hand.

\paragraph{Succinct certificates and light clients.} Mina~\cite{mina}
came first with a constant-size recursive certificate attesting the whole
history; Molt Petit's is useful only with the recency check, giving
agreement on recent chains. Plumo~\cite{plumo} syncs ultralight clients
with SNARK-proved committee hand-offs, argued on paper.
NiPoPoWs~\cite{nipopow} and FlyClient~\cite{flyclient} certify
proof-of-work weight probabilistically; Molt Petit proves a deterministic
validity rule, with no sampling.

\paragraph{BFT with certain finality.} Like Tendermint~\cite{tendermint}
and HotStuff~\cite{hotstuff}, Molt Petit has a fixed roster, rotating
leaders and certain finality, but exchanges no votes: the same $2n/3$
threshold is enforced by counting blocks in a window, and the light
client's object is a proof of that rule, not a bundle of votes.

\paragraph{Long-range attacks.} Keys that leak after the fact are the
classic long-range threat~\cite{stakebleeding}. Key-evolving
signatures~\cite{forwardsec}, as in Praos, make an old key useless once
erased; weak subjectivity~\cite{weaksubj} has a long-offline client fetch
a trusted checkpoint. Molt Petit's recency check is weak subjectivity
turned into a theorem assumption, checked against the client's own clock
rather than a trusted party, and the windowed budget says exactly which
leaks a fork can still use: only those recency rejects.

\paragraph{Machine-checked consensus.} Velisarios~\cite{velisarios},
Verdi~\cite{verdi} and IronFleet~\cite{ironfleet} prove running consensus
code correct by extracting it from a model inside the prover. Molt Petit
proves a vote-free density consensus with a recursive light-client
certificate, and \emph{imports} the deployed sources into
Lean~\cite{lean,mathlib} with Charon + Aeneas~\cite{aeneas} and
Thales~\cite{thales}.

\section{Limitations and future work}
\label{sec:limitations}

\paragraph{What the design gives up.} The roster is fixed and closed
(keys rotate, validators do not change). Validators have equal weight and
a public schedule: a predictable leader is easier to censor, but no
randomness beacon is needed. There is no recovery from halts: if a
complete window ever has fewer than $\quorum$ blocks, the chain halts for
good; at $n = 7$ and $\tau \approx 4$\,s, three validators missing their
slots within one window of about $28$\,s do it. Safety reads a clock
(Assumption~\ref{ass:clock}), but never a synchronous network.

\paragraph{What is not machine-checked.} Beyond the trusted base of
Section~\ref{sec:impl}, every unchecked fact is a named assumption: the
deployment assumptions of Section~\ref{sec:assumptions}; honest custody per key for mode 3 and the untimed
mode 2; version independence under theft; the roster's commitment to the
cold roots, the circuit's trust anchor; that the clock reading handed
to the client is the node's real clock; the mode-3 duties of moving together and erasure; and that a lost key
exposes nothing, which rests on a signing record that survives the loss.

\paragraph{The long-range residual.} Exposure after a window's stretch
does not count against it, so a patient attacker can collect late
signatures into an arbitrarily long fork on old slots. Recency rejects
it, and Theorem~\ref{thm:forge} forbids preparing one in advance, but it
is never impossible. Forward-secure signatures~\cite{forwardsec} with erasure, or
counting every exposure forever,
would remove it, at the price of a stateful signer or of charging a
cleaned-up signer for good. Key rotation does not. We keep the windowed
budget and accept the contained residual.

\paragraph{Artifact.} \code{lake build} re-checks every theorem cited
here, with axiom guards (the \code{Molt} library's executable functions
are proved equal to the model's, \code{reference\_eq\_model}).
\code{nix flake check} re-checks that the vendored translations are what
the pinned translators produce, and runs the \code{rust-keyrot/} tests. 

\begin{thebibliography}{99}

\bibitem{ouroboros}
A.~Kiayias, A.~Russell, B.~David, R.~Oliynykov.
\emph{Ouroboros: A Provably Secure Proof-of-Stake Blockchain Protocol}.
CRYPTO 2017.

\bibitem{praos}
B.~David, P.~Gazi, A.~Kiayias, A.~Russell.
\emph{Ouroboros Praos: An Adaptively-Secure, Semi-synchronous
Proof-of-Stake Blockchain}. EUROCRYPT 2018.

\bibitem{samasika}
J.~Bonneau, I.~Meckler, V.~Rao, E.~Shapiro.
\emph{Ouroboros Samasika: A Proof-of-Stake Consensus Protocol for
Succinct Blockchains}. O(1)~Labs / Mina Protocol technical paper, 2020.

\bibitem{genesis}
C.~Badertscher, P.~Gazi, A.~Kiayias, A.~Russell, V.~Zikas.
\emph{Ouroboros Genesis: Composable Proof-of-Stake Blockchains with
Dynamic Availability}. ACM CCS 2018.

\bibitem{weaksubj}
V.~Buterin. \emph{Proof of Stake: How I Learned to Love Weak
Subjectivity}. Ethereum blog, 2014.

\bibitem{mina}
J.~Bonneau, I.~Meckler, V.~Rao, E.~Shapiro.
\emph{Mina: Decentralized Cryptocurrency at Scale}. Whitepaper, 2020.
IACR ePrint 2020/352.

\bibitem{plumo}
P.~Vesely, K.~G\"urkan, M.~Straka, A.~Gabizon, P.~Jovanovic,
G.~Konstantopoulos, A.~Oines, M.~Olszewski, E.~Tromer.
\emph{Plumo: An Ultralight Blockchain Client}. Financial Cryptography
2022.

\bibitem{nipopow}
A.~Kiayias, A.~Miller, D.~Zindros.
\emph{Non-Interactive Proofs of Proof-of-Work}. Financial Cryptography
2020.

\bibitem{flyclient}
B.~B\"unz, L.~Kiffer, L.~Luu, M.~Zamani.
\emph{FlyClient: Super-Light Clients for Cryptocurrencies}. IEEE S\&P
2020.

\bibitem{tendermint}
E.~Buchman, J.~Kwon, Z.~Milosevic.
\emph{The latest gossip on BFT consensus}. arXiv:1807.04938, 2018.

\bibitem{hotstuff}
M.~Yin, D.~Malkhi, M.~K.~Reiter, G.~Golan~Gueta, I.~Abraham.
\emph{HotStuff: BFT Consensus in the Lens of Blockchain}. PODC 2019.

\bibitem{canettisig}
R.~Canetti.
\emph{Universally Composable Signature, Certification, and
Authentication}. CSFW 2004.

\bibitem{forwardsec}
M.~Bellare, S.~K.~Miner.
\emph{A Forward-Secure Digital Signature Scheme}. CRYPTO 1999.

\bibitem{keyinsulated}
Y.~Dodis, J.~Katz, S.~Xu, M.~Yung.
\emph{Key-Insulated Public Key Cryptosystems}. EUROCRYPT 2002.

\bibitem{keyinsulatedsig}
Y.~Dodis, J.~Katz, S.~Xu, M.~Yung.
\emph{Strong Key-Insulated Signature Schemes}. PKC 2003.

\bibitem{stakebleeding}
P.~Ga\v{z}i, A.~Kiayias, A.~Russell.
\emph{Stake-Bleeding Attacks on Proof-of-Stake Blockchains}. Crypto Valley
Conference on Blockchain Technology (CVCBT) 2018.

\bibitem{sphincsplus}
D.~J.~Bernstein, A.~H\"ulsing, S.~K\"olbl, R.~Niederhagen, J.~Rijneveld,
P.~Schwabe.
\emph{The SPHINCS+ Signature Framework}. ACM CCS 2019.

\bibitem{fips205}
National Institute of Standards and Technology.
\emph{Stateless Hash-Based Digital Signature Standard}. FIPS 205, 2024.

\bibitem{xmssrfc}
A.~H\"ulsing, D.~Butin, S.~Gazdag, J.~Rijneveld, A.~Mohaisen.
\emph{XMSS: eXtended Merkle Signature Scheme}. RFC 8391, 2018.

\bibitem{poseidon}
L.~Grassi, D.~Khovratovich, C.~Rechberger, A.~Roy, M.~Schofnegger.
\emph{Poseidon: A New Hash Function for Zero-Knowledge Proof Systems}.
USENIX Security 2021.

\bibitem{velisarios}
V.~Rahli, I.~Vukoti\'c, M.~V\"olp, P.~Esteves-Verissimo.
\emph{Velisarios: Byzantine Fault-Tolerant Protocols Powered by Coq}.
ESOP 2018.

\bibitem{verdi}
J.~R.~Wilcox, D.~Woos, P.~Panchekha, Z.~Tatlock, X.~Wang, M.~D.~Ernst,
T.~Anderson.
\emph{Verdi: A Framework for Implementing and Formally Verifying
Distributed Systems}. PLDI 2015.

\bibitem{ironfleet}
C.~Hawblitzel, J.~Howell, M.~Kapritsos, J.~R.~Lorch, B.~Parno,
M.~L.~Roberts, S.~Setty, B.~Zill.
\emph{IronFleet: Proving Practical Distributed Systems Correct}. SOSP
2015.

\bibitem{lean}
L.~de~Moura, S.~Ullrich.
\emph{The Lean 4 Theorem Prover and Programming Language}. CADE 2021.

\bibitem{mathlib}
The mathlib Community.
\emph{The Lean Mathematical Library}. CPP 2020.

\bibitem{thales}
J.~Alama. \emph{Thales: a TypeScript-to-Lean compiler}.
\url{https://github.com/jessealama/thales}.

\bibitem{aeneas}
S.~Ho, J.~Protzenko. \emph{Aeneas: Rust Verification by Functional
Translation}. ICFP 2022. (With the Charon frontend,
\url{https://github.com/AeneasVerif}.)

\bibitem{plonky2}
Polygon Zero. \emph{Plonky2: Fast Recursive Arguments with PLONK and
FRI}. \url{https://github.com/0xPolygonZero/plonky2}.

\end{thebibliography}

\end{document}
